\chapter{Alpha-Thalassemia Genetic Testing}

\section*{What Is This Test?}
Alpha-thalassemia genetic testing identifies deletions or sequence variants in the alpha-globin genes, usually HBA1 and HBA2. It is useful because routine hemoglobin electrophoresis can be normal or only subtly abnormal in alpha-thalassemia trait.

\section*{Alternative Names}
Alpha-Globin Gene Testing; HBA1/HBA2 Testing; Alpha-Thalassemia Molecular Analysis.

\section*{Why This Test Is Ordered}
Testing is used for persistent microcytosis with normal iron studies, suspected hemoglobin H disease, unexplained anemia, newborn or prenatal evaluation, and carrier or partner testing when reproductive risk is relevant.

\section*{How This Test Is Performed}
DNA is usually extracted from blood and tested for common alpha-globin deletions, followed by sequencing or deletion/duplication analysis when the initial panel is negative or the clinical suspicion remains high.

\section*{Preparation, Risks, and Limitations}
No fasting is required. Risks are limited to blood collection. A limited deletion panel may miss uncommon variants. Results should be interpreted with CBC indices, iron studies, hemoglobin analysis, family history, and ethnicity or ancestry when relevant.

\section*{Understanding Results}
One-gene loss usually represents a silent carrier state, while two-gene loss causes alpha-thalassemia trait. Three-gene loss can cause hemoglobin H disease, and four-gene loss causes severe alpha-thalassemia with fetal or neonatal consequences. Precise parental genotypes are important for reproductive counseling.

\section*{References}
Thalassaemia International Federation guidance on alpha-thalassemia diagnosis and carrier assessment.

\clearpage
\section*{Understanding Results and Follow-up}
The laboratory report should identify the specimen, method, units, reference interval or decision limit, and whether the result is positive, negative, abnormal, or inconclusive. Reference intervals vary with age, sex, pregnancy, specimen type, assay, and laboratory; a result outside a range is not automatically a diagnosis, and a result within a range does not always exclude disease. Discuss unexpected results with the ordering clinician, who can decide whether repeat or confirmatory testing, treatment, or referral is appropriate. Do not change medicines or delay urgent care based on an isolated result.


\section*{Regulatory Status and Authorization Date}
This chapter describes a test category, not a specific commercial product. FDA, CDSCO, EU IVDR, and MHRA approval or registration dates therefore cannot be assigned without the assay manufacturer, product name, instrument, and intended use. For laboratory-developed tests, an individual agency approval date may not exist. See the Preface for the interpretation of regulatory status.

\clearpage
